% ============================================================================
%  附录 B：Python 量化工具链
% ============================================================================
\chapter{Python 量化工具链}
\label{app:python}

本附录给出从数据到回测的最小可用工具链。代码全部为可直接运行的 Python 片段，
变量与注释使用英文以保证跨平台可执行。

\section{环境与依赖}

\begin{lstlisting}[style=shellstyle,
  caption={推荐的环境准备方式（uv 与 conda 二选一）},
  label={lst:env}]
# Option 1: uv (fast, modern)
uv venv .venv && source .venv/bin/activate
uv pip install numpy pandas scipy statsmodels arch scikit-learn cvxpy \
               matplotlib pyarrow yfinance numba

# Option 2: conda
conda create -n quant python=3.11 -y
conda activate quant
pip install -r requirements.txt
\end{lstlisting}

\begin{table}[htbp]
\centering
\caption{量化金融常用 Python 库与分工}\label{tab:python-libs}
\tablefont
\begin{tabularx}{\textwidth}{@{}L{3.2cm}Y L{4.6cm}@{}}
\toprule
\textbf{库} & \textbf{职责} & \textbf{典型用途} \\
\midrule
\texttt{numpy} & 向量化数值计算、随机数生成 & 蒙特卡洛路径、线性代数 \\
\texttt{pandas} & 带标签的时间序列与表格运算 & 数据对齐、重采样、滚动窗口 \\
\texttt{scipy} & 优化、插值、统计分布、求根 & 隐含波动率反解、曲线拟合 \\
\texttt{statsmodels} & 计量经济学模型 & OLS/GLS、ARIMA、协整、面板 \\
\texttt{arch} & 条件异方差模型 & GARCH/EGARCH/GJR、波动率预测 \\
\texttt{scikit-learn} & 机器学习与预处理 & 因子非线性建模、交叉验证 \\
\texttt{cvxpy} & 凸优化建模（DSL） & 均值--方差、风险平价、CVaR 优化 \\
\texttt{numba} & 即时编译加速 & 逐路径回测、期权定价循环 \\
\texttt{pyarrow / parquet} & 列式存储 & 大规模行情与因子数据 \\
\bottomrule
\end{tabularx}
\end{table}

\section{数据读取与清洗}

\begin{lstlisting}[style=pythonstyle,
  caption={价格数据对齐、复权与缺失值处理},
  label={lst:data}]
import numpy as np
import pandas as pd

def load_prices(path: str, tickers: list[str]) -> pd.DataFrame:
    """Load adjusted close prices and align them on a common calendar."""
    px = pd.read_parquet(path)                 # MultiIndex or wide format
    px = px.loc[:, tickers].sort_index()
    px = px[~px.index.duplicated(keep="last")]
    px = px.ffill(limit=5)                     # short gaps only
    px = px.dropna(how="any")                  # keep the common history
    return px

def log_returns(px: pd.DataFrame) -> pd.DataFrame:
    r = np.log(px / px.shift(1)).dropna(how="all")
    return r.replace([np.inf, -np.inf], np.nan).dropna(how="any")
\end{lstlisting}

\begin{warningbox}[三个最常见的清洗错误]
\begin{enumerate}[leftmargin=1.6em]
  \item \textbf{用未复权价格计算收益}：除权除息日会出现虚假的巨额负收益，因子测试结果被彻底污染；
  \item \textbf{用 \texttt{ffill()} 无长度限制地填补停牌与退市}：会把「无交易」伪装成「零波动」，
        压低协方差估计并制造虚假的平稳性；
  \item \textbf{只保留当前仍在交易的股票}：典型的幸存者偏差，对价值与困境反转类因子的影响尤其大。
\end{enumerate}
\end{warningbox}

\section{描述统计与肥尾检验}

\begin{lstlisting}[style=pythonstyle,
  caption={收益率描述统计：年化、偏度、峰度、Jarque--Bera},
  label={lst:desc}]
import numpy as np
import pandas as pd
from scipy import stats

def summary_stats(r: pd.Series, periods: int = 252) -> pd.Series:
    ann_ret = r.mean() * periods
    ann_vol = r.std(ddof=1) * np.sqrt(periods)
    downside = r[r < 0].std(ddof=1) * np.sqrt(periods)
    return pd.Series({
        "ann_return": ann_ret,
        "ann_vol": ann_vol,
        "sharpe": ann_ret / ann_vol if ann_vol > 0 else np.nan,
        "sortino": ann_ret / downside if downside > 0 else np.nan,
        "skew": stats.skew(r),
        "excess_kurtosis": stats.kurtosis(r),
        "jb_pvalue": stats.jarque_bera(r).pvalue,
        "max_drawdown": (r.add(1).cumprod().cummax()
                         - r.add(1).cumprod()).max(),
    })
\end{lstlisting}

\begin{keybox}[判断「肥尾」的正确姿势]
\texttt{excess\_kurtosis} 为正只说明样本呈现出厚尾，\textbf{并不证明总体四阶矩存在}。
在 $t$ 分布且 $\nu\le 4$ 时，样本峰度的抽样分布极不稳定。更稳健的做法是：
(1) 用 Hill 估计或 GPD 拟合尾指数；(2) 直接比较经验分位数与正态分位数在 1\%/0.1\% 分位上的差异。
\end{keybox}

\section{波动率建模}

\begin{lstlisting}[style=pythonstyle,
  caption={用 arch 拟合 GARCH(1,1)-t 并预测年化波动率},
  label={lst:garch}]
import numpy as np
from arch import arch_model

def garch_forecast(returns: np.ndarray, horizon: int = 5) -> dict:
    """Fit GARCH(1,1) with Student-t errors and forecast volatility."""
    r = 100.0 * returns                       # percentage scale is stabler
    res = arch_model(r, mean="Constant", vol="Garch", p=1, q=1,
                     dist="t").fit(disp="off")
    fc = res.forecast(horizon=horizon, reindex=False)
    var = fc.variance.values[-1]              # in (%)^2
    ann_vol = np.sqrt(var * 252.0) / 100.0    # back to decimals
    alpha, beta = res.params["alpha[1]"], res.params["beta[1]"]
    half_life = np.log(0.5) / np.log(alpha + beta)
    return {"params": res.params, "ann_vol": ann_vol,
            "half_life_days": half_life, "loglik": res.loglikelihood}
\end{lstlisting}

\begin{notebox}[持久性与半衰期]
GARCH(1,1) 的波动率冲击半衰期为 $-\ln 2/\ln(\alpha+\beta)$。当 $\alpha+\beta=0.95$ 时约 13.5 天，
$0.99$ 时约 69 天。若估计出的 $\alpha+\beta$ 逼近甚至超过 1，说明模型认为波动率冲击永不衰减，
此时点预测虽然可用，但长期方差 $\omega/(1-\alpha-\beta)$ 已无意义——这通常意味着样本中存在
结构突变或长记忆，应考虑 FIGARCH 或加入机制转换。
\end{notebox}

\section{因子研究}

\begin{lstlisting}[style=pythonstyle,
  caption={因子中性化与信息系数（IC）计算},
  label={lst:factor}]
import numpy as np
import pandas as pd
import statsmodels.api as sm

def neutralize(factor: pd.Series, exposures: pd.DataFrame) -> pd.Series:
    """Cross-sectional OLS: factor ~ industry dummies + style exposures."""
    y = factor.dropna()
    X = exposures.loc[y.index]
    X = sm.add_constant(X)
    resid = sm.OLS(y, X, hasconst=True).fit().resid
    return resid

def rank_ic(factor: pd.DataFrame, fwd_ret: pd.DataFrame) -> pd.Series:
    """Cross-sectional Spearman correlation per date."""
    ic = {}
    for dt in factor.index.intersection(fwd_ret.index):
        f, r = factor.loc[dt], fwd_ret.loc[dt]
        mask = f.notna() & r.notna()
        if mask.sum() > 30:
            ic[dt] = f[mask].corr(r[mask], method="spearman")
    return pd.Series(ic).sort_index()

def ic_report(ic: pd.Series, periods: int = 252) -> dict:
    return {"IC_mean": ic.mean(), "IC_std": ic.std(ddof=1),
            "ICIR": ic.mean() / ic.std(ddof=1) * np.sqrt(periods),
            "t_stat": ic.mean() / ic.std(ddof=1) * np.sqrt(len(ic)),
            "positive_ratio": (ic > 0).mean()}
\end{lstlisting}

\begin{appbox}[IC 与 IR 的经验尺度]
横截面 RankIC 的月频均值若能达到 $0.03$--$0.05$ 已属良好，$0.05$ 以上通常意味着
较强的选股能力（或尚未被修正的前视偏差）。月频 ICIR 在 $0.5$ 左右是可用的水准。
注意：\textbf{IC 的绝对值远小于直觉预期}，这正是「基本定律」中需要用广度 $\sqrt{\text{Breadth}}$
去放大信号的原因。
\end{appbox}

\section{组合优化}

\begin{lstlisting}[style=pythonstyle,
  caption={用 cvxpy 求解带约束的均值--方差与等风险贡献组合},
  label={lst:opt}]
import numpy as np
import cvxpy as cp

def min_variance(mu: np.ndarray, cov: np.ndarray, w_max: float = 0.1):
    n = len(mu)
    w = cp.Variable(n)
    cons = [cp.sum(w) == 1, w >= 0, w <= w_max]
    prob = cp.Problem(cp.Minimize(cp.quad_form(w, cov)), cons)
    prob.solve()
    return w.value

def max_utility(mu: np.ndarray, cov: np.ndarray, lam: float = 5.0):
    n = len(mu)
    w = cp.Variable(n)
    cons = [cp.sum(w) == 1, w >= 0]
    obj = cp.Maximize(mu @ w - 0.5 * lam * cp.quad_form(w, cov))
    cp.Problem(obj, cons).solve()
    return w.value

def risk_parity(cov: np.ndarray, budget: np.ndarray | None = None):
    """Solve for equal (or budgeted) risk contributions via convex program."""
    n = cov.shape[0]
    b = np.ones(n) / n if budget is None else budget / budget.sum()
    y = cp.Variable(n, nonneg=True)
    port_var = cp.quad_form(y, cov)
    obj = cp.Minimize(port_var - b @ cp.log(y))
    cp.Problem(obj).solve()
    w = y.value / y.value.sum()
    return w
\end{lstlisting}

\begin{intuitionbox}[风险平价的建模技巧]
上面把「等风险贡献」写成了凸优化问题，而不是直接解非线性方程组。技巧在于：
最小化 $y^T\Sigma y - \sum_i b_i\ln y_i$ 的一阶条件恰好是
$2(\Sigma y)_i = b_i/y_i$，即 $y_i(\Sigma y)_i \propto b_i$，
还原成权重后就是「风险贡献等于预算」。这类\textbf{把均衡条件写成优化目标}的思路，
在组合管理中非常常用。
\end{intuitionbox}

\section{衍生品定价}

\begin{lstlisting}[style=pythonstyle,
  caption={Black--Scholes 定价、Newton 法反解隐含波动率与蒙特卡洛验证},
  label={lst:bsm}]
import numpy as np
from scipy.stats import norm

def bsm_price(S, K, r, sigma, T, kind="call"):
    d1 = (np.log(S / K) + (r + 0.5 * sigma ** 2) * T) / (sigma * np.sqrt(T))
    d2 = d1 - sigma * np.sqrt(T)
    if kind == "call":
        return S * norm.cdf(d1) - K * np.exp(-r * T) * norm.cdf(d2)
    return K * np.exp(-r * T) * norm.cdf(-d2) - S * norm.cdf(-d1)

def implied_vol(price, S, K, r, T, kind="call", tol=1e-8, max_iter=100):
    sigma = 0.2
    for _ in range(max_iter):
        diff = bsm_price(S, K, r, sigma, T, kind) - price
        vega = S * norm.pdf((np.log(S / K) + (r + 0.5 * sigma ** 2) * T)
                            / (sigma * np.sqrt(T))) * np.sqrt(T)
        if abs(diff) < tol:
            break
        sigma -= diff / vega
    return sigma

def bsm_mc(S, K, r, sigma, T, n=200_000, seed=42):
    rng = np.random.default_rng(seed)
    z = rng.standard_normal(n)
    ST = S * np.exp((r - 0.5 * sigma ** 2) * T + sigma * np.sqrt(T) * z)
    payoff = np.maximum(ST - K, 0.0)
    return np.exp(-r * T) * payoff.mean(), payoff.std(ddof=1) / np.sqrt(n)
\end{lstlisting}

\begin{warningbox}[Newton 法反解隐含波动率的失效场景]
当 Vega 趋于 0 时（深度虚值、临近到期），Newton 迭代会剧烈震荡或发散。
实务中的处理方式是：(1) 先做无套利检查（价格是否落在内在价值与标的价之间）；
(2) 用二分法或 Brent 法作为兜底；(3) 对波动率曲面做 SVI/SABR 参数化后再插值，
而不是逐点反解。对全期限、全行权价的逐点反解很容易产生「日历套利」与「蝶式套利」的假象。
\end{warningbox}

\section{工程实践要点}

\begin{table}[htbp]
\centering
\caption{从研究到生产的常见工程陷阱}\label{tab:eng-pitfalls}
\tablefont
\begin{tabularx}{\textwidth}{@{}L{3.4cm}Y@{}}
\toprule
\textbf{陷阱} & \textbf{症状与对策} \\
\midrule
随机数不可复现 &
回测结果无法重现。对策：统一使用 \texttt{np.random.default\_rng(seed)}，并把种子写入配置。\\
浮点比较 &
用 \texttt{==} 比较价格或权重导致误判。对策：用 \texttt{np.isclose} 与绝对/相对容差。\\
循环写得太慢 &
Python 逐行循环处理数百万行行情。对策：向量化，或对嵌套路径循环使用 \texttt{numba.njit}。\\
内存爆炸 &
一次性加载十年全市场逐笔数据。对策：Parquet + 按日期分区、分块处理、\texttt{dtype} 降位。\\
时间对齐错误 &
把不同时区的收盘价直接相减，或用未来公布的财报。对策：统一时区与交易日历，严格使用
「公告日 + 1」作为可用日。\\
参数硬编码 &
换一个市场就全盘失效。对策：把交易成本、涨跌停、停牌规则做成显式参数并记录版本。\\
\bottomrule
\end{tabularx}
\end{table}

\begin{keybox}[可复现性是量化研究的第一生产力]
一份研究结果如果不能被别人（包括三个月后的你自己）用同一份数据与代码复现出来，
它的价值为零。最小要求是：\textbf{数据快照 + 代码版本 + 随机种子 + 参数配置}四件套齐全，
并且把「研究环境」与「生产环境」的差异显式记录在案。
\end{keybox}
